\section*{Introdução}
\addcontentsline{toc}{section}{Introdução}
\label{sec:introducao}

% (==>) Este é o fenômeno; esta é a questão; há evidência científica suficiente para justificar investigá-la

% ECONOMIA + INSTITUIÇÕES --> DECISÃO DE USO DA TERRA --> ESTOQUE AMBIENTAL
% Foco do Capítulo 1 (H1): como retornos econômicos esperados e o ambiente institucional condicionam a conversão da vegetação nativa.
% (A retroalimentação — vegetação × clima --> produtividade — constitui o objeto de H2 no Capítulo 2; este artigo investiga a formação e transformação do estoque).

% HIPÓTESE DA TESE INVESTIGADA NESTE ARTIGO (H1 e desdobramentos):
% ---
% H1: RETORNOS, INSTITUIÇÕES E CONVERSÃO
% Retorno econômico esperado + ambiente institucional --> decisão de uso da terra --> taxa de conversão de vegetação nativa
%
% H1-a: PREÇOS E CAPACIDADE PRODUTIVA LOCAL
% Choques favoráveis de preços agrícolas geram maior pressão de conversão onde a capacidade produtiva da terra é mais elevada.
%
% H1-b: CONDICIONAMENTO INSTITUCIONAL
% Regras, monitoramento, fiscalização e sanções condicionam a sensibilidade da taxa de conversão ao retorno econômico esperado.
%
% Heterogeneidade: respostas distintas entre atividades (lavoura vs. pastagem) e entre biomas ( Amazônia, e Pantanal).

% EIXOS CONCEITUAIS QUE CONECTAM O ARTIGO À TESE:
% ---
% ENDOGENEIDADE AMBIENTAL
% O meio ambiente é tratado como componente endógeno do processo econômico: seu estoque é transformado pelas decisões produtivas sob incentivos e instituições.
%
% GOVERNANÇA TERRITORIAL
% Instituições, regulação e mecanismos de implementação operam sobre agentes e territórios heterogêneos, condicionando a intensidade da resposta aos incentivos de mercado.
%
% INTERTEMPORALIDADE
% Decisões presentes de conversão de uso da terra alteram estoques ambientais acumulados e condicionam oportunidades e riscos econômicos futuros.

% ------------------------------------
%%% (I) CONTEXTO
% ------------------------------------
% Mapeamento bibliográfico de H1:
% - Dilema conservação vs. rentabilidade e determinantes da decisão privada: Angelsen (2010), Lubowski et al. (2008), Jung et al. (2022).
% - Implementação regulatória e enforcement local: Viallon et al. (2019), Assunção et al. (2023), Pacheco & Meyer (2022).
% - Preços relativos, custos e contexto territorial: Harding et al. (2021), Chomitz & Thomas (2003).
% - Reorganização produtiva vs. expansão de fronteira: Bragança (2018).
% - Mato Grosso como contexto empírico e diversidade de trajetórias: VanWey et al. (2013), Simões et al. (2020).

%% Conflito: preservar ou converter (dilema econômico e determinantes privados/institucionais da decisão).
A expansão agropecuária costuma ser analisada à luz do dilema entre conservar a vegetação nativa, promovendo a sustentabilidade ambiental, ou converter a terra para explorar usos potencialmente mais rentáveis. Esse dilema, contudo, não esgota os fatores que condicionam a decisão de conversão. Além da aptidão agronômica e da rentabilidade esperada da produção, a expectativa de valorização da terra também compõe o cálculo dos agentes econômicos ao buscarem otimizar o uso de suas propriedades \cite{PachecoMeyer2022}. Trata-se, afinal, de uma decisão privada moldada por preços, custos, restrições de crédito, direitos de propriedade, regras ambientais e seus mecanismos de implementação \cite{Angelsen2010}. É nesse contexto amplo que a qualidade do solo e os retornos de usos alternativos se tornam determinantes para explicar as escolhas dos proprietários \cite{LubowskiEtAl2008}. No cenário brasileiro, essa dinâmica ganha contornos institucionais específicos: os efeitos do CAR \footnote{O Cadastro Ambiental Rural (CAR), instituído pelo art. 29 da Lei nº 12.651/2012, é um registro público eletrônico obrigatório para todos os imóveis rurais do Brasil, destinado ao monitoramento ambiental e ao combate ao desmatamento. Proprietários ou possuidores declaram informações georreferenciadas sobre o imóvel, a vegetação nativa, as Áreas de Preservação Permanente e a Reserva Legal, sujeitas à análise e validação pelo órgão ambiental competente.} sobre a conversão para pastagens, por exemplo, variam conforme o porte da propriedade e seu histórico de conformidade, operando possivelmente por meio de alterações no acesso ao crédito \cite{JungEtAl2022}.

%% Heterogeneidade institucional e implementação local (regras formais vs. efetividade da fiscalização/enforcement).
Isto posto, uma mesma regra formal pode produzir respostas diferentes conforme as condições locais de implementação. Como exemplos, a análise de arranjos regulatórios na Suíça mostrou a relevância do enforcement estatal e dos direitos e deveres legais dos proprietários \cite{ViallonEtAl2019}, ao tempo em que, igualmente, na Amazônia brasileira há evidências de que monitoramento e fiscalização efetivos podem de fato ajudar a conter o desmatamento de forma eficaz \cite{AssuncaoEtAl2023}. Sob essa perspectiva, a conversão ou conservação da vegetação deve ser tratada como resultado endógeno da interação entre os retornos econômicos dos usos alternativos da terra e os incentivos, custos e riscos associados ao ambiente institucional \cite{PachecoMeyer2022}.

% Ponte para H1: retornos econômicos, incentivos institucionais e evolução do estoque de vegetação nativa.
Nessa toada, mudanças nos preços relativos, na rentabilidade agrícola, preços fundiários ou nas condições de acesso aos mercados podem aumentar o retorno esperado da conversão, enquanto fiscalização, sanções e restrições legais podem alterar a resposta dos agentes a esses incentivos \cite{HardingEtAl2021}. Nçao menos importante, a intensidade dessa resposta também pode variar conforme o contexto econômico e as características territoriais do bioma ou região \cite{ChomitzThomas2003}. Ato contínuo, compreender a evolução do estoque de vegetação requer analisar os retornos econômicos associados ao uso da terra e as instituições sob as quais esses retornos são percebidos e transformados em decisões de conversão ou conservação.

% Nuance teórica e empírica: expansão de fronteira vs. reorganização em áreas já convertidas (preços relativos e conservação espontânea/egoísta).
Porém, essa relação entre estoque de vegetação e expectativas econômicas, ressalte-se, não implica necessariamente que, por exemplo, uma elevação dos preços dos produtos agropecuários conduza ao aumento do desmatamento e de outros danos ambientais. Exemplificando, na bacia do Tapajós \footnote{A quinta maior sub-bacia da Amazônia brasileira, com uma área de aproximadamente 764 mil quilômetros quadrados, fica localizada nos estados de Mato Grosso e do Pará, com pequenas porções nos estados do Amazonas e de Rondônia.}, \citeonline{Braganca2018} encontra que a elevação do preço relativo de lavouras em relação à carne favoreceu a substituição de pastagens por cultivos (\textit{pasture-to-cropland}) e reduziu o desmatamento local. A pressão para abrir novas áreas agropastoris deve, assim, ser distinguida da reorganização da produção em áreas já convertidas. 

Assim sendo, essa distinção também importa para a avaliação das políticas de conservação, pois seus incentivos e restrições podem modificar tanto a pressão para abrir novas áreas quanto as escolhas produtivas em áreas já convertidas. A classificação de políticas e mecanismos de conservação como instrumentos regulatórios ou econômicos, por si só, não esclarecem como os produtores respondem a essas intervenções. Dentre outros exemplos, as obrigações de conservação previstas no Código Florestal \cite{SticklerEtAl2013}, os critérios de compra estabelecidos pela Moratória da Soja \footnote{A Moratória da Soja (2006) é o primeiro acordo setorial privado de desmatamento zero implementado nos trópicos. O mecanismo proíbe as principais indústrias processadoras e comerciais de grãos de comprarem soja cultivada em áreas desmatadas do bioma Amazônia após julho de 2008. Formalizado de maneira independente da legislação estatal e validado pela suprema corte do Brasil em 2026 face às pressões das bancadas agrícolas regionais, o pacto opera como uma barreira não tarifária de mercado baseada em monitoramento geoespacial privado. Suas implicações teóricas e empíricas são centrais para a economia ambiental, pois demonstram a eficácia de sanções comerciais privadas na mitigação de externalidades ambientais, ao mesmo tempo em que introduzem assimetrias econômicas, reconfiguram o preço da terra e geram vazamento de desmatamento (\textit{spillover effect}) para biomas vizinhos sem restrições semelhantes.} \cite{HardingEtAl2021} e as receitas potenciais associadas à conservação do carbono florestal \cite{Angelsen2010,SticklerEtAl2013} atuam por mecanismos distintos sobre as decisões de uso da terra. Deste modo, a resposta dos produtores depende de como essas intervenções modificam os retornos dos usos alternativos da terra, das possibilidades de ajuste da produção e das condições locais de implementação. É necessário, portanto, verificar empiricamente se a implementação dessas intervenções reduz a conversão da vegetação nativa \cite{Angelsen2010}.

%% Mato Grosso como caso empírico: grande escala agropecuária, três biomas e diferenciais biofísicos e econômicos.
Nesse contexto, Mato Grosso constitui um caso empírico especialmente informativo para investigar como os retornos econômicos dos usos agropecuários e o ambiente institucional se relacionam com as decisões de conversão da vegetação nativa. O estado combina produção agropecuária em larga escala com três biomas -- Amazônia, Cerrado e Pantanal --, cujas peculiaridades ambientais e econômicas resultam em diferenças relevantes de vegetação, clima, disponibilidade hídrica, sistemas produtivos e custo de oportunidade da terra dentro do mesmo território \cite{VanWeyEtAl2013}. Por esse norte, a literatura documenta tanto a expansão e a intensificação agrícola quanto a diversidade das trajetórias de uso da terra em Mato Grosso \cite{SimoesEtAl2020}.

%%%%% MUDAR (INÍCIO) --> 
% Em H1 eu lido com a questão da governança territorial e a heterogeneidade espacial. 
% Em H2 eu lido com a questão da intertemporalidade e a questão da heterogeneidade espacial.
%% Heterogeneidade territorial como mecanismo econômico: variação das respostas entre Amazônia, e Pantanal.
Desse modo, a heterogeneidade espacial, ambiental e econômica observada em Mato Grosso permite examinar se os retornos econômicos e as instituições produzem respostas distintas de conversão entre atividades e biomas. Tudo considerado, não basta identificar se uma regra ambiental ou uma alteração na rentabilidade afeta a conversão em termos médios, importa verificar se a magnitude dessas respostas depende das condições econômicas e territoriais sobre as quais os incentivos incidem. Assim, a diversidade entre Amazônia, Cerrado e Pantanal interessa à pesquisa como parte integrante dos próprios mecanismos econômicos em análise, pois as restrições biofísicas e os marcos regulatórios definem custos de oportunidade da terra e potenciais de rentabilidade distintos para as atividades agropecuárias.
%%%%% MUDAR (FIM) <--

% ------------------------------------
%%% (II) PROBLEMA, OBJETIVOS E JUSTIFICATIVA
% ------------------------------------

%% Problema central e formulação da questão de pesquisa (H1 em Mato Grosso).
Diante desse quadro, o problema central deste artigo consiste em compreender como incentivos econômicos e institucionais moldam a conversão ou manutenção da vegetação em ambientes ecologicamente e economicamente heterogêneos. O \textbf{problema} pode ser sintetizado no seguinte questionamento:

\medskip
\noindent\emph{Como retornos econômicos e instituições condicionam as decisões de conversão ou manutenção da vegetação nativa em Mato Grosso, e em que medida essa relação varia conforme a capacidade produtiva local, as atividades agropecuárias e os biomas?}
\medskip

%% Desdobramento do problema nas duas dimensões empíricas de H1: (i) preços × capacidade produtiva; (ii) condicionamento institucional.
Essa pergunta desdobra-se em duas dimensões empiricamente relacionadas: (i) se mudanças nos preços dos produtos agropecuários se traduzem em maior conversão onde a capacidade produtiva para esses usos é mais elevada; e (ii) se as condições institucionais modificam a intensidade dessa resposta. Em conjunto, essas dimensões buscam esclarecer como incentivos privados e instituições interagem na relação entre produção e conservação, oferecendo subsídios para o desenho e a focalização territorial de políticas ambientais.

%% Objetivo geral e objetivos específicos (i)–(iv): conversão bruta, resposta a preços × capacidade, instituições e heterogeneidade.
O \textbf{objetivo geral} deste artigo é investigar como e em que intensidade os determinantes observáveis do retorno econômico esperado e o ambiente institucional se relacionam com as decisões de conversão da vegetação nativa em Mato Grosso. De forma específica, o artigo busca: (i) construir medidas de conversão bruta para agricultura e pastagem; (ii) examinar a resposta dessas medidas a variações de preços combinadas com capacidade produtiva predeterminada; (iii) avaliar o condicionamento institucional dessa resposta, distinguindo associação de identificação causal; e (iv) verificar sua heterogeneidade entre atividades e biomas.

%%%%% MUDAR (INÍCIO) --> 
% Em H1 eu lido com a questão da governança territorial e a heterogeneidade espacial. 
% Em H2 eu lido com a questão da intertemporalidade e a questão da heterogeneidade espacial.
%% Relevância teórica e eixos da tese: meio ambiente endógeno, dimensão territorial e intertemporal; divisão analítica com o Capítulo 2 (H2).
A relevância deste trabalho decorre, em termos teóricos, de tratar o \textbf{meio ambiente} como componente \textbf{endógeno} do processo econômico, cujo estoque é transformado por decisões de produção e pelas instituições que as condicionam. Essa relação envolve uma \textbf{dimensão territorial}, na medida em que incentivos e instituições podem gerar respostas distintas conforme as condições locais, e uma \textbf{dimensão intertemporal}, pois a conversão presente altera estoques ambientais e oportunidades econômicas futuras. No âmbito da tese, este artigo examina a formação desse estoque pela decisão de uso da terra; o capítulo 2, dedicado à H2, examina se a vegetação previamente existente modifica a resposta da produtividade agrícola a choques climáticos. A articulação dos capítulos permitirá discutir a contribuição produtiva do meio ambiente, sem atribuir a este primeiro artigo um teste do circuito completo de retroalimentação. A condição de insumo produtivo é compatível com a existência de serviços ambientais cujos benefícios ultrapassam a propriedade e geram externalidades.
%%%%% MUDAR (FIM) <--

%% Relevância aplicada e contribuição para políticas: decisões privadas, fiscalização pública, governança e limites institucionais.
Em termos empíricos e aplicados, compreender essas relações pode orientar decisões de produtores sobre conversão, conservação e manejo da terra e subsidiar governos e organizações responsáveis por fiscalização e alocação de recursos ambientais. A contribuição pretendida está em revelar quando condições institucionais e territoriais alteram a resposta aos incentivos econômicos. Isso não pressupõe a ineficácia dos instrumentos existentes nem a superioridade automática da governança local: a literatura admite diferentes combinações de regras, monitoramento e sanções, cuja efetividade depende do contexto \cite{AnderssonEtAl2014}.

% ------------------------------------
%%% (III) HIPÓTESES
% ------------------------------------

%% H1 — Governança territorial e heterogeneidade espacial: investiga como incentivos econômicos e instituições condicionam as decisões de conversão ou preservação da vegetação, considerando diferenças entre territórios e biomas.

%% Enunciação da hipótese H1 da tese e de seus desdobramentos empíricos (H1-a e H1-b).
Dito isso, o artigo examina a primeira hipótese da tese e seus dois desdobramentos:

\medskip
% Hipótese geral do artigo (H1):
\noindent\textbf{(H1) Retornos, instituições e conversão.}

% CLAIM: H1-C001
% Maiores retornos econômicos esperados dos usos agropecuários elevam a taxa de conversão de vegetação nativa.
\noindent\emph{A taxa de conversão de vegetação nativa para usos agropecuários aumenta quando se eleva o retorno econômico esperado desses usos.}

\medskip
% Desdobramento empírico H1-a: interação entre choques de preços e capacidade produtiva predeterminada.
\noindent\textbf{(H1-a) Preços e capacidade produtiva local.}

\noindent\emph{O efeito de aumentos nos preços dos produtos agropecuários sobre a taxa de conversão de vegetação nativa é maior onde a capacidade produtiva local da terra para esses usos é mais elevada.}

\medskip
% Desdobramento empírico H1-b: condicionamento institucional sobre a sensibilidade da conversão aos retornos.
\noindent\textbf{(H1-b) Condicionamento institucional.}

\noindent\emph{A resposta da taxa de conversão ao retorno econômico esperado dos usos agropecuários é condicionada pelo ambiente institucional, particularmente pelas regras, pelo monitoramento, pela fiscalização e pelas sanções que alteram os custos e os riscos associados à conversão.}
\medskip

%% Delimitação conceitual e variável dependente: vegetação nativa, conversão agropecuária e taxa de conversão.
Nesta hipótese, vegetação nativa corresponde à cobertura vegetal natural existente no território, enquanto sua conversão designa a substituição dessa cobertura por usos agropecuários, como lavouras ou pastagens. A taxa de conversão refere-se à parcela da vegetação nativa disponível em determinado território e período que é convertida para esses usos, constituindo o principal resultado de interesse. A ausência de conversão observada tampouco equivale, necessariamente, a uma decisão ativa de conservação ou à integridade ecológica da vegetação remanescente.

%% Identificação do retorno esperado: aproximação por determinantes observáveis (preços e capacidade produtiva como implicação observável de H1).
O retorno econômico esperado dos usos agropecuários corresponde ao benefício privado esperado da utilização produtiva da terra, considerando receitas e custos associados à atividade. Empiricamente, esse retorno não é diretamente observado nas bases reunidas, podendo ser aproximado por seus determinantes econômicos observáveis. Entre eles, destacam-se os preços dos produtos agropecuários e a capacidade produtiva local da terra. A H1-a expressa a expectativa de que choques favoráveis de preços produzam maior pressão sobre a conversão onde a terra oferece maior capacidade produtiva para os usos beneficiados. O teste dessa previsão será, portanto, um teste de uma implicação observável de H1, condicionado à validade da aproximação; não uma estimativa direta do retorno líquido esperado.

%% Ambiente institucional: regras formais e mecanismos de implementação (monitoramento, fiscalização e sanções).
O ambiente institucional compreende as regras formais e os mecanismos de implementação que alteram os incentivos econômicos associados à conversão, incluindo direitos de propriedade, regulação ambiental, monitoramento, fiscalização e sanções. A H1-b não pressupõe que essas instituições determinem isoladamente a conversão, mas que modifiquem a intensidade com que alterações nos retornos econômicos esperados se traduzem em conversão efetiva. Quando uma medida representar, de maneira defensável, maior efetividade da fiscalização, espera-se menor sensibilidade da conversão ao retorno. Contagens de autos ou valores de multas, por si sós, não satisfazem essa condição.

%% Heterogeneidade empírica (atividades e biomas), escopo da análise e critérios para resultados nulos ou contrários.
As relações previstas em H1, H1-a e H1-b podem apresentar heterogeneidade entre atividades agropecuárias e biomas. Diferentes atividades apresentam tecnologias, estruturas de custos e retornos distintos, enquanto os biomas diferem em condições biofísicas e produtivas. Essa heterogeneidade será tratada como dimensão empírica da análise, sem constituir uma hipótese autônoma. Resultados nulos, de sinal contrário ou restritos a determinados grupos serão discutidos como limites ou evidência contrária às previsões propostas.

% ------------------------------------
%%% (IV) MÉTODO E ORGANIZAÇÃO DO ARTIGO
% ------------------------------------

%% Estratégia empírica: dados em painel (2001–2024), fontes, ordenação temporal e identificação econométrica.
As hipóteses serão examinadas com dados em painel para Mato Grosso, tendo 2001--2024 como período candidato, conforme a delimitação já adotada na preparação da qualificação. A análise combinará transições anuais de uso da terra, preços agropecuários, medidas predeterminadas de capacidade produtiva e registros institucionais. As especificações considerarão a ordenação temporal entre incentivos e conversão, efeitos fixos territoriais e temporais e controles compatíveis com a identificação. A endogeneidade da fiscalização, a dependência espacial e a validade temporal das medidas de capacidade produtiva receberão tratamento explícito. A amostra final depende da harmonização territorial e da obtenção das conversões brutas, ainda pendentes.

%% Organização do artigo: estrutura das seções subsequentes (1 a 5).
Além desta introdução, o artigo está organizado em cinco seções. A primeira revisa a literatura sobre retornos, instituições e conversão da vegetação. A segunda apresenta Mato Grosso como contexto empírico, descreve os dados e explicita a construção das variáveis. A terceira desenvolve a estratégia econométrica para examinar as hipóteses do artigo. A quarta estabelece os critérios de apresentação e discussão dos resultados a serem obtidos, com atenção à heterogeneidade e à robustez. Por fim, a quinta reúne as contribuições pretendidas e as limitações do estudo. Nesta versão para qualificação, não são apresentados resultados econométricos.

